\section{Discussion}

\label{sec:discussion}

\subsection{Training Support}

Candle currently focuses on inference. Autograd support exists but is not optimized for large-scale training. Training support is planned as future work, building on the inference engine's memory efficiency.

\subsection{Limitations}

\begin{enumerate}[leftmargin=1.4em]
    \item \textbf{Ecosystem maturity}: The Rust ML ecosystem has fewer pre-trained models and utilities compared to Python.
    \item \textbf{Learning curve}: Rust's ownership system has a steeper learning curve than Python for ML practitioners.
    \item \textbf{Dynamic shapes}: Some operations with fully dynamic shapes cannot benefit from compile-time type checking.
    \item \textbf{AMD GPU}: ROCm support is experimental; CUDA remains the primary GPU backend.
\end{enumerate}

\subsection{Future Work}

\begin{itemize}[leftmargin=1.1em]
    \item \textbf{Distributed inference}: Tensor parallelism and pipeline parallelism for multi-GPU serving.
    \item \textbf{WebAssembly}: Compile models to WASM for browser-based inference.
    \item \textbf{Training optimization}: Gradient checkpointing and mixed-precision training support.
    \item \textbf{Speculative decoding}: Implement Medusa and draft-verify decoding for faster generation.
\end{itemize}

